\title{Large Language Models Can Automatically Engineer Features for Few-Shot Tabular Learning}

\begin{abstract}
Large Language Models (LLMs) have demonstrated exceptional capabilities in addressing complex and novel reasoning tasks, making them highly promising for tabular learning, which is fundamental to numerous real-world scenarios. This paper introduces an innovative in-context learning approach, \textsf{FeatLLM}, that leverages LLMs as feature constructors to generate a dataset specifically designed for optimal tabular prediction. The features produced by \textsf{FeatLLM} facilitate class probability estimation using a basic downstream machine learning algorithm, such as linear regression, resulting in robust few-shot learning performance. The framework is designed to rely solely on this elementary predictive model utilizing the derived features during inference. Unlike traditional LLM-powered solutions, \textsf{FeatLLM} removes the necessity of querying the LLM for every individual instance at inference. Furthermore, it functions with only API-level interaction with LLMs, thereby sidestepping prompt size constraints. Empirical results on diverse tabular datasets from various domains confirm that \textsf{FeatLLM} consistently generates superior quality rules and achieves substantially better results (on average a 10\% improvement) over methods like TabLLM and STUNT.
\end{abstract}

\section{Introduction}
Recent advancements have shown that Large Language Models (LLMs) excel at producing responses for unseen tasks without the need for extensive, task-specific adaptation~\cite{creswell2022selection,imani2023mathprompter,taylor2022galactica}. The broad textual datasets these models are trained on endow them with substantial prior expertise that can substitute domain experts across numerous fields~\cite{Genomes2021,singhal2023large,wilcox2003role}, enabling automation in various applications such as dialogue analysis~\cite{finch2023leveraging} and software maintenance~\cite{fan2023automated}. Notably, providing LLMs with task context and several illustrative examples through prompts enables them to discern task relevance and focus attention on salient features and samples~\cite{brown2020language,zhao2023survey}, suggesting their aptitude for sophisticated data interpretation and feature engineering.

The application of LLMs' prior knowledge and reasoning strengths has expanded into tabular learning. For example, understanding that advanced age correlates with higher disease risk can enhance disease forecasting. Current approaches articulate tasks and feature details in natural language, serialize tabular data, and input these into LLMs~\cite{dinh2022lift,wang2023anypredict}. Subsequently, methods such as in-context learning~\cite{nam2023semi} or efficient fine-tuning~\cite{dinh2022lift,hegselmann2023tabllm} are often applied. The integration of LLM-informed prior knowledge has repeatedly shown greater effectiveness than traditional tabular models, especially under limited data circumstances~\cite{hegselmann2023tabllm}.

Existing approaches for LLM-based tabular data analysis encounter several notable challenges. For full end-to-end prediction workflows, each individual data instance requires a separate LLM evaluation, contributing to substantial computational overhead. Achieving high predictive accuracy typically demands fine-tuning the LLM, which becomes impractical for LLMs whose training infrastructure or fine-tuning APIs are unavailable—this issue is especially critical as the latest high-performing LLMs limit access to inference-only APIs~\cite{anil2023palm,openai2023gpt4}. Moreover, these methods are frequently hindered by prompt length constraints~\cite{liu2023lost}; as the tabular data's feature set expands, the serialized prompt exceeds feasible limits, often resulting in inoperable models or degraded performance. Collectively, these issues pose significant obstacles for deploying LLM-powered tabular models in practice.

In contrast, our method reimagines the role of LLMs within tabular prediction tasks by leveraging them for feature engineering rather than direct end-to-end prediction. This strategy utilizes the extensive prior knowledge encapsulated in LLMs more effectively. We propose a new in-context learning framework, \textsf{FeatLLM}, designed to uncover the ``criteria'' embedded in the LLM's decision-making process. For example, in the context of disease classification, the LLM can deduce and articulate explicit rules indicating which feature patterns correspond to disease detection. These rules—or criteria—then form new features that replace the original ones for downstream analysis. By shifting the LLM's function to generating task-specific features, subsequent predictive models can be streamlined, thus reducing inference times and obviating the need for sophisticated models once informative features are in place. Feature extraction is facilitated by in-context learning, requiring only prompt-based example demonstrations, thereby circumventing traditional model training. Additionally, ensemble-inspired approaches~\cite{breiman2001random}, such as feature bagging, help address prompt size limitations, ensuring scalability even with extensive feature spaces.

\textsf{FeatLLM} organizes the prompt for feature engineering into two distinct sub-tasks: (1) comprehending the underlying problem and deducing how features interact with targets; and (2) utilizing insights from the first sub-task along with a limited set of examples to formulate discriminative rules for class separation. This sequential reasoning helps LLMs disregard irrelevant features during rule development. The constructed rules are subsequently applied to the dataset (interpreted as program code), converting samples into binary indicators that represent compliance with each rule. Ultimately, a simple model is used to predict class probabilities based on these newly derived features. The framework's resilience is enhanced through an ensemble mechanism that repeatedly creates features while employing feature bagging, which strategically manages the number of features and samples per batch, alleviating issues related to large prompt sizes. \textsf{FeatLLM} operates independently of traditional training requirements and delivers efficient inference, broadening the applicability of LLMs to a wide spectrum of practical tasks.

In our experiments, \textsf{FeatLLM} was tested on 13 tabular datasets under few-shot conditions, demonstrating strong and consistent performance. We illustrate that LLMs are capable of leveraging both external domain expertise and knowledge extracted from the data, leading to the creation of informative features. The proposed framework surpasses state-of-the-art few-shot learning baselines in diverse scenarios. The source code is made available through an anonymized GitHub repository at~\texttt{https://github.com/Sungwon-Han/FeatLLM}.

\section{Related Work}

\subsection{Few-Shot Learning with Tabular Data}

Few-shot learning has made it possible to acquire generalizable data representations while significantly lowering annotation expenses~\cite{chen2018closer}. Although initial developments focused primarily on visual data, this paradigm has since shown strong generalizability across multiple formats, including natural language, audio signals, and—more recently—tabular datasets~\cite{majumder2022few,nam2023stunt,schick2021exploiting}. The challenge of extracting meaningful information from limited annotated samples often increases the susceptibility to spurious patterns~\cite{bartlett2021deep}. To counteract this, researchers have sought to leverage external sources or embed domain-driven priors, effectively encouraging the models toward productive inductive biases. Illustrative strategies include harnessing large pools of unlabeled tabular samples combined with intelligent data augmentation methods for semi-supervised approaches~\cite{ucar2021subtab,yoon2020vime}, or performing task-independent unsupervised pre-training to achieve robust initializations~\cite{somepalli2021saint}. Such unsupervised pre-training objectives commonly utilize masking, noise injection, or data corruption, followed by reconstructive learning~\cite{arik2021tabnet,majmundar2022met}, and alternatively, deploying data augmentation for contrastive loss formulations~\cite{bahri2022scarf,chen2020simple}. Nevertheless, the intrinsic diversity present in tabular structures poses difficulties in crafting augmentations that generalize well, and these techniques have proven to be limited within few-shot scenarios~\cite{nam2023stunt}.

Emerging methods propose training models across a wide array of unrelated tabular datasets, then employing transfer learning mechanisms to adapt to the downstream tasks~\cite{zhang2023generative,levin2022transfer,zhu2023xtab}. For example, TransTab~\cite{wang2022transtab} applies transformer models to infer semantic connections among columns using column names, supplementing value-based learning. This type of transfer of relational information enables successful zero-shot or few-shot evaluations on unseen tables. TabPFN~\cite{hollmann2022tabpfn}, on the other hand, creates synthetic mixtures of data distributions for pre-training Bayesian neural architectures. When inference is needed, both the training and testing data from the specific task are input to the model without further optimization. The broad prior amassed from pre-training supports performance improvements over classic tree-based models and confers advantages in low-data regimes.

\subsection{Language-Interfaced Tabular Learning}

Integrating prior knowledge into tabular learning has been further advanced by the application of large language models (LLMs)~\cite{anil2023palm,openai2023gpt4}. Due to their extensive training over diverse text collections, LLMs possess strong generalization capabilities, enabling them to address novel problems in several domains~\cite{brown2020language}. Recent developments leverage this by encoding tabular information into textual format and incorporating it into the input prompts of LLMs~\cite{dinh2022lift,hegselmann2023tabllm,wang2023anypredict}. Presenting tabular values, accompanying feature details, and specific task instructions as textual prompts allows the models to deduce complex task-feature relationships for downstream prediction. Progress in this domain has included both specialized training with efficient parameter adaptation techniques such as LoRA~\cite{hu2021lora} and IA3~\cite{liu2022few}, as well as in-context learning approaches where few-shot demonstrations are provided to the LLMs in prompts, enabling reasoning and inference without additional model updates~\cite{dinh2022lift,hegselmann2023tabllm,nam2023semi,slack2023tablet}.

While previous approaches have shown success in advancing few-shot tabular learning, their reliance on routing inference through the LLM for every prediction can hinder scalability and practicality within industrial settings. This work proposes a shift away from the typical end-to-end, black-box strategy where each input's output is determined by passing it through an LLM. Rather, it focuses on utilizing the LLM to explicitly derive conditions or decision rules corresponding to each class label. \textsf{FeatLLM} transforms these generated rules into programmatic code, producing novel features that allow subsequent predictions to occur independently of the LLM. The method makes use of in-context learning to formulate these rules by harnessing prior knowledge and leveraging information gleaned from the few-shot data samples.

\section{Methods}
The central goal of \textsf{FeatLLM} is to identify "rules," or specific conditions related to each possible answer class, from a small set of example inputs, as illustrated in Figure~\ref{fig:main_model}. These rules are extracted by the LLM and subsequently translated into binary programmatic features for each sample, denoting whether the rule criteria are met. The resulting binary features are used in a dot product computation with non-negative trainable weights, producing estimates of class probabilities. Training this model enables learning which features are most influential in determining class assignment (see Section~\ref{Sec:training}). This process is iteratively executed with bagging, and the individual outcomes are aggregated through ensembling. The subsequent sections detail the problem setup and describe the procedures for each phase.

\vspace*{-0.12in}
\paragraph{Problem formulation.} Consider a tabular dataset $\mathcal{D} = \{(\mathbf{x}^i, \mathbf{y}^i)\}_{i=1}^{N}$ containing $N$ labeled instances. Each $\mathbf{x}^i$ is a vector of length $d$ corresponding to $d$ input features, where feature names $F = \{f_j\}_{j=1}^{d}$ are expressed in natural language and paired with their respective definitions. Under a classification framework, each label $\mathbf{y}^i$ represents one class from a predetermined set. In $k$-shot learning scenarios, the model is trained on a randomly selected subset of $k$ ($<N$) labeled examples.

\subsection{Prompt Design for Extracting Rules}
\label{Sec:prompt_design}
In order to facilitate rule extraction by the LLM using a more faithful reasoning trajectory, we craft the prompting strategy to reflect a process akin to an expert's methodology in tabular machine learning. The prompt is structured into three primary sections (illustrated in Fig.~\ref{fig:prompt_for_rule} and further detailed in Appendix~\ref{sec:example_rule}):

\vspace*{-0.12in}
\paragraph{Basic information description.} This segment delivers key information necessary for addressing the task (highlighted in orange in Fig.~\ref{fig:prompt_for_rule}). It comprises the task statement (which includes labeling guidance), explanations of feature attributes, and demonstration examples. The task statement is framed as an inquiry (for instance, ``Does this patient's myocardial infarction complications data indicate chronic heart failure? Yes or no?"; see more in Appendix~\ref{sec:task_desc}). Each feature is described by specifying its value type and, when available, its definition. For categorical variables, the range of potential category values may be provided. Demonstration examples consist of a few training cases converted to textual format together with their corresponding true labels. This serialization method adheres to conventions used in earlier work~\cite{dinh2022lift,hegselmann2023tabllm}:

\begin{align}
&\text{Serialize}(\mathbf{x}^i,\mathbf{y}^i, F) = \nonumber \\ 
&``\ f_1\ \text{is}\ \mathbf{x}^i_1.\ ... \ f_d\  \text{is}\  \mathbf{x}^i_d.\ \ \text{Answer:}\  \mathbf{y}^i\ ”
\end{align}

\vspace*{-0.12in}
\paragraph{Reasoning instruction.} Our objective is to strengthen the reasoning abilities of the LLM by furnishing it with explicit reasoning instructions (highlighted in blue in Fig.~\ref{fig:prompt_for_rule}). The instruction starts with an opening sentence that mirrors the chain-of-thought paradigm~\cite{wei2022chain}. Subsequently, we segment the inference process into two distinct phases. During the initial phase, the LLM is prompted to deduce causal associations or tendencies between feature variables and the given task description, independent of any example demonstrations. This approach relies upon general or commonsense knowledge, ensuring the LLM leverages its prior information about the scenario at hand. In the subsequent phase, the model incorporates example demonstrations, as well as insights from the previous stage, to infer rules corresponding to each class. This bifurcated reasoning method helps avoid the identification of irrelevant or spurious relationships, thus directing attention toward more pertinent features. To strike a balance between insufficient and excessive rule extraction—which could either cause underfitting or unnecessarily inflate the prompt length—a fixed rule count is established (namely, 10)\footnote{A detailed discussion concerning rule count is presented in Section~\ref{sec:analysis}.} per class. Furthermore, when constructing the rules, we instruct the LLM to reference the feature type found in the basic information description to ensure the compatibility of feature types within the formulated rules.

\vspace*{-0.12in}
\paragraph{Response instruction.} For more effective parsing and utilization of the rules inferred by the LLM, we provide directives specifying the desired structure of the model's responses (displayed in yellow in Fig.~\ref{fig:prompt_for_rule}). The prompt outlines both the required response template for each answer class (i.e., response format) and the prescribed structure for individual rules (i.e., rule format). Such instructions help the LLM choose the correct structure according to each feature's type. In scenarios without supplemental guidance, the LLM is permitted to synthesize multiple rules using logical connectives, such as AND or OR. An exemplary result illustrating this can be found in Appendix~\ref{sec:outcome-rule}.

\subsection{Inferring Class Likelihood via Rules}
\label{Sec:training}

\paragraph{Rule parsing for automated feature creation.} The rules extracted during the earlier phase are employed to construct new binary features. For each class, these features reflect whether a given sample complies with its respective rule set—represented as either '0' or '1'. Such features facilitate assessment of a sample's propensity for belonging to specific classes. Nevertheless, since the LLM yields rules in natural language format, an effective parsing mechanism is essential for generating these features programmatically. Although rule generation is regulated with prescribed response and rule formatting to simplify parsing, the LLM's outputs can still exhibit unpredictable syntactic variation~\cite{kovavc2023large}. As an example, the LLM may replace symbols such as ``$>$” and ``$<$” with written expressions like ``greater than" or ``less than," complicating automated parsing.

Rather than developing sophisticated parsing algorithms to handle these text irregularities, we capitalize on the LLM’s own capabilities. The LLM is adept at grasping semantic content even with syntactic inconsistencies and can synthesize concise code snippets when prompted (due to its exposure to code datasets). To take advantage of this, we frame prompts by including function names, detailed descriptions for inputs and outputs, and the extracted rules, which are then supplied to the LLM. Illustrative prompts and their results are presented in Figures~\ref{fig:prompt_for_parsing} and ~\ref{sec:example_parsing}-\ref{sec:example_parse_outcome} in the Appendix. Subsequently, the code produced is run via Python’s \texttt{exec()} method to achieve the necessary data transformation.

\vspace*{-0.12in}
\paragraph{Class likelihood estimation.} With the new rule-derived binary features per class, a straightforward approach to estimate the likelihood of a sample’s class membership is to tally the number of satisfied rules for each class (that is, sum up the relevant binary features for a class). Nonetheless, individual rules and features may contribute unevenly, prompting the need to derive their importance based on training data. To address this, we employ a standard machine learning (ML) model, applied separately for the binary features of each class, to infer their respective weights. For instance, a bias-free linear model can be used. Denote the binary feature vector constructed from sample $\mathbf{x}^i$ for class $k$ as $\mathbf{z}_k^i \in \{0, 1\}^R$, with $R$ indicating the rule count for each class and $c$ as the total number of classes. The probability of each class $\mathbf{p}^i$ is determined as follows:

\begin{align}
\text{logit}_k^i &= \max(\mathbf{w}_k, 0) \cdot \mathbf{z}_k^i, \nonumber \\
\mathbf{p}^i &= \text{Softmax}({[\text{logit}_{1}^i, ..., \text{logit}_{c}^i]}). \label{eq:infer_prob}
\end{align}

In this context, $\mathbf{w}_k$ refers to the learnable weights within the linear model, which denote the significance of each input feature. By restricting these weights to reside in the positive subspace, it is possible to guarantee that the features produced by the LLM only augment, rather than diminish, the likelihood associated with each class during training. Subsequently, the logit vector is transformed into class probabilities using the softmax function. The optimization objective is the cross-entropy loss, with $\mathbf{w}_k$ updated based on limited, few-shot examples.

\vspace*{-0.12in}
\paragraph{Ensembling with bagging.} The entire pipeline is executed repeatedly, yielding a suite of models that together constitute an ensemble for final predictions. Let the collection of parameters from $T$ runs be $\{\mathbf{w}[t]\}_{t=1}^T$, each producing a probability vector $\mathbf{p}^i$ via $\mathbf{w}[t]$, and the ensemble prediction is calculated by taking the mean over these probabilities in accordance with Eq.~\ref{eq:infer_prob}.

To promote variability in rule formation for the ensemble, the LLM’s temperature setting is increased during inference. Additionally, the few-shot exemplars are randomly permuted, capitalizing on evidence that different ordering may affect LLM outputs~\cite{lu2022fantastically}\footnote{See Appendix~\ref{sec:diversity} for discussion on rule diversity}. Further leveraging ensemble diversity, bagging is applied so that each iteration utilizes a different subset of features or data samples; this strategy becomes particularly advantageous as the number of samples or features grows. The ensemble procedure brings two primary benefits. Firstly, should erroneous rules emerge in some runs, accurate rules generated elsewhere can offset them, improving robustness. This phenomenon aligns with the principle of LLM self-consistency~\cite{wang2022self}, since rules that repeatedly appear across independent trials are more apt to be correct. Secondly, the ensemble addresses restrictions posed by prompt length limitations—when the tabular dataset surpasses prompt capacity, bagging enables selective sampling to preserve model reliability. Even though an individual rule set may lack comprehensive feature or instance coverage, combining the outputs of multiple weak learners generated from different trials alleviates these limitations, thereby strengthening overall predictive reliability.

\section{Experiments}
We assess the effectiveness of \textsf{FeatLLM} across several tabular datasets and carry out a range of analyses to explore the behavior and characteristics of our framework. Owing to limitations in manuscript length, supplementary experiments—including those involving alternative LLMs, qualitative assessments, and additional insights—are presented in the Appendix.

\subsection{Performance Evaluation}
\paragraph{Datasets.}
Our study leverages 13 datasets spanning binary and multi-class classification problems: (1) Adult~\cite{asuncion2007uci} tasks the model with predicting if an individual's income surpasses \$50,000 per year; (2) Bank~\cite{moro2014data} concerns forecasting a customer’s decision regarding term deposit subscription; (3) Blood~\cite{yeh2009knowledge} predicts whether blood donors will return for future donations; (4) Car~\cite{kadra2021well} assesses car quality; (5) Communities~\cite{misc_communities_and_crime_183} estimates crime rates in various localities; (6) Credit-g~\cite{kadra2021well} determines if an individual constitutes a favorable or unfavorable credit risk; (7) Diabetes\footnote{\texttt{kaggle.com/datasets/uciml/pima-indians-diabetes-database}} classifies the likelihood of diabetes in patients; (8) Heart\footnote{\texttt{kaggle.com/datasets/fedesoriano/heart-failure-prediction}} targets detection of coronary artery disease; and (9) Myocardial~\cite{misc_myocardial_infarction_complications_579} identifies individuals suffering from chronic heart failure.

Moreover, we incorporate newer and synthetic datasets that are seldom examined in prior literature and were not present during LLM pre-training: (10) Cultivars, which aims to predict the grain yield level for soybean cultivars\footnote{\texttt{archive.ics.uci.edu/dataset/913}}; (11) NHANES, designed for classification of age groups from person records\footnote{\texttt{archive.ics.uci.edu/dataset/887}}; (12) Sequence-type, which sorts sequences into four possible categories: arithmetic, geometric, Fibonacci, or Collatz; (13) Solution-mix, which evaluates whether the concentration of a mixture comprising four solutions exceeds 0.5. The first two datasets were only made available in the UCI Machine Learning Repository following September 2023, which guarantees their absence from the pre-training data of the LLM API (gpt-3.5-turbo-0613) utilized by our model. The last two represent synthetic data, precluding any possibility of prior memorization by the LLM API and ensuring a fair test scenario.

These datasets exhibit considerable variability in terms of sample count and attribute complexity, with further specifics outlined in Table~\ref{Tab:basic-desc}. Each dataset includes explicit names and attribute descriptions, while those such as `Communities' and `Myocardial' feature over a hundred attributes, presenting difficulties attributable to the restricted context window of LLMs.

\vspace*{-0.12in}
\paragraph{Baselines.}
We benchmark our model against nine baseline methods. The initial trio of baselines reflects standard techniques for tabular data: (1) Logistic regression (LogReg), (2) XGBoost~\cite{chen2016xgboost}, and (3) RandomForest~\cite{ho1995random}. Following this, two baselines operate under the assumption of access to large volumes of unlabeled data: (4) SCARF~\cite{bahri2022scarf} and (5) STUNT~\cite{nam2023stunt}. It is important to note that gathering extensive unlabeled data is often impractical in genuine operational contexts. Among the recent contenders, (6) TabPFN~\cite{hollmann2022tabpfn} synthesizes varied datasets for model pretraining. The last group consists of three LLM-based methods: (7) In-context learning (In-context)~\cite{wei2022emergent}, (8) TABLET~\cite{slack2023tablet}, and (9) TabLLM~\cite{hegselmann2023tabllm}. In-context learning places several training instances within the input prompt directly, foregoing any parameter updates. TABLET enriches prompts using additional clues extracted from an external classifier, employing rule-sets and prototypes for more accurate inference. TabLLM adopts a parameter-efficient strategy—specifically IA3~\cite{liu2022few}—to fine-tune the LLM using a limited set of samples.

\vspace*{-0.12in}
\paragraph{Implementation details.}
\textsf{FeatLLM} is built upon GPT-3.5 as its primary LLM backbone, but its architecture remains compatible with alternative LLMs (refer to Appendix~\ref{sec:backbones} for experimental outcomes involving varying backbones). Inference parameters for the LLM are configured with a temperature of 0.5 and top-p set to 1, adhering to API defaults. Ensemble size is set at 20, while rule extraction employs 10 rules. Insights into the effects of hyper-parameter choices can be found in Figure~\ref{fig:hyperparameter2} and Appendix~\ref{sec:hyper-parameter}.
The linear model training utilizes the Adam optimizer at a learning rate of 0.01 across 200 epochs. Epoch selection is conducted through $k$-fold cross-validation to ensure optimality. Baseline reproduction uses GPT-3.5 for in-context learning; for TabLLM, T0~\cite{sanh2022multitask} is adopted per its original configuration. It's important to highlight that TabLLM limits LLM selection to those with openly released checkpoints, making GPT-3.5 unavailable for TabLLM experimentation. In applying conventional machine learning algorithms (LogReg, XGBoost, RandomForest), parameters are tuned by combining grid search with $k$-fold cross-validation. The chosen $k$ is either 2 or 4, ensuring each training partition contains at least one sample from every class. More comprehensive implementation specifics are provided in Appendix~\ref{sec:baselines}.

\vspace*{-0.12in}
\paragraph{Results.}
Tables~\ref{tab:main1} and \ref{tab:main2} summarize comparative results spanning various datasets. Each experimental setting is executed three times with different random splits of training and test data, and we report the mean and standard deviation of the resulting AUC scores. Across these benchmarks, \textsf{FeatLLM} either emerges as the highest-performing method or places second, relative to other baseline models. Remarkably, our approach matches the performance of traditional tabular baselines employing 32 shots, even when restricted to only 4 shots (refer to Fig.~\ref{fig:compares}). It is also worth mentioning that, although STUNT and TabPFN attained notable improvements on a subset of datasets, their aggregate gains over standard machine learning methods remained modest. In addition, LLM-based methods including in-context learning, TABLET, and TabLLM were unable to carry out inference on tabular datasets with high-dimensional feature spaces. In contrast, our framework, which leverages both bagging and ensemble strategies, remains robust and yields satisfactory results even when dealing with datasets comprising a large number of features. Notably, our bagging and ensemble framework could be integrated into existing LLM-powered models; however, such integration would require twice as many inferences per sample, significantly escalating computational demands and thereby rendering it impractical.

\subsection{Analysis \& Discussion}
\label{sec:analysis}
\paragraph{Ablation study.} To assess the individual impact of key components, we conducted ablation studies, considering the following variations: (1) \textsf{-Tuning}: the weight tuning step in the linear model is skipped, and instead, logits are obtained by directly summing the binary features; (2) \textsf{-Ensemble}: the ensemble procedure is excluded; (3) \textsf{-Description}: the use of feature descriptions is omitted; and (4) \textsf{-Reasoning}: Step-1 of the reasoning instruction for rule generation is left out.

In Table~\ref{tab:ablation}, we present the mean changes in AUC resulting from each ablation, averaged across all evaluated datasets. Without exception, omitting any of the tested components led to diminished performance. Of particular note, the value added by weight tuning increases in scenarios with more available shots, suggesting that with a larger dataset, tuning enables more reliable identification of significant rules and has a correspondingly larger effect. Conversely, in few-shot settings, providing feature descriptions and reasoning instructions yields more substantial improvements, highlighting the significance of leveraging LLMs' prior knowledge for boosting performance when training data is limited. Finally, the use of ensembles consistently confers performance improvements under all tested conditions.

\vspace*{-0.12in}
\paragraph{Training \& inference time comparison.} We assess the training and inference durations across several methods, conducting experiments using the Adult dataset with a training sample size of 16 shots. The majority of approaches utilized a single A100 GPU; however, TabLLM operated on four A100 GPUs due to model parallelism requirements. Table~\ref{tab:cost} reports these time comparisons. \textsf{FeatLLM} demonstrates particularly efficient inference time, matching the performance of standard methods. By contrast, other LLM-centric techniques, including In-context learning, TABLET, and TabLLM, incur notably higher inference times. With regard to training duration, \textsf{FeatLLM} aligns with other LLM-based solutions. It necessitates only 30 API queries—this is not costly in terms of resources and can readily be parallelized to minimize time overhead.

\vspace*{-0.12in}
\paragraph{Quality of features generated by \textsf{FeatLLM}.}
To examine the effectiveness of rule-based features produced by \textsf{FeatLLM} for real-world tasks, we carry out an experiment that quantifies the informativeness of these features. Concretely, the average AUROC between the newly constructed features and the target labels is calculated, and for each dataset, we determine the proportion of features that achieve an AUROC exceeding 0.5. The outcomes are presented in Table~\ref{tab:quality}. These results substantiate that most rule-based features are pertinent to the target variable and supply valuable information for the associated task (see ``Before tuning'' column). Furthermore, during linear model training, features deemed less relevant—those with weights below a specified threshold—are systematically excluded from the final model (see ``After tuning'' column). The threshold is set at 0.1, informed by the overall distribution of learned weights.

\vspace*{-0.12in}
\paragraph{Effect of adding spurious correlations.} We evaluated how well different approaches mitigate the impact of spurious correlations under conditions where labeled examples are scarce. For this analysis, we employed the Heart dataset, systematically appending random columns from the Adult dataset that have no bearing on the predictive objectives of the Heart dataset. By incrementally increasing the number of these extraneous columns, we tracked changes in model performance. To ensure an equitable evaluation, all methods were limited to 8 labeled samples and the absence of unlabeled data, thus excluding techniques such as SCARF or STUNT.

Figure~\ref{fig:spurious} illustrates model performance as additional spurious correlations are introduced. Among the tested methods, \textsf{FeatLLM} maintained performance most effectively, exhibiting minimal deterioration even as the volume of irrelevant data grew. This resilience likely stems from its explicit reliance on prior knowledge when establishing task-feature associations during rule extraction. This approach curtails the formulation of rules involving non-informative features, enhancing robustness. Empirically, only 1–2\% of generated rules utilized noisy columns. However, our results also indicate that \textsf{FeatLLM} is susceptible to learning unintended correlations and misinterpreting causal relationships if columns from datasets related to the target task are appended at random (refer to Appendix~\ref{sec:spurious-appendix} for additional discussion).

\vspace*{-0.12in}
\paragraph{Hyper-parameter analysis}
\textsf{FeatLLM} relies principally on two hyper-parameters: the ensemble count and the number of rules generated per class for each LLM query. To understand their influence on overall performance, we conducted empirical assessments. Our findings indicate that increasing the ensemble size leads to enhanced performance; however, after reaching approximately 20 ensembles, the benefits plateau (refer to Fig.~\ref{fig:appendix-hyperparameter1} in the Appendix). With respect to rule quantity, no statistically meaningful variation in performance was observed across different values. Nevertheless, we selected 10 rules as the optimal setting, balancing prompt complexity against prediction quality (see Fig.~\ref{fig:hyperparameter2}). We hypothesize that while a greater number of rules contributes more information by expanding the input space, it may also introduce overfitting—especially under low-shot circumstances.

\section{Conclusion}
This work introduces \textsf{FeatLLM}, which elevates the standard for LLM-driven few-shot tabular data modeling. Unlike traditional approaches that leverage LLMs for direct sample-wise predictions, \textsf{FeatLLM} emphasizes criteria generation akin to feature engineering. By employing rule induction alongside a bagging-based ensemble, the method substantially reduces inference expenses, operates exclusively via API without necessitating training, and avoids typical limitations associated with feature dimensionality. Notably, \textsf{FeatLLM} demonstrates robust predictive accuracy and offers an effective, data-efficient solution for deploying LLMs on real-world tabular datasets.

Designed specifically for low-shot learning scenarios, \textsf{FeatLLM}'s extension plans involve broadening its ability to devise novel features in larger datasets, investigating feature types beyond rules, and enhancing interpretability. These future developments aim to generalize its application to a wider variety of tasks.

\section{Broader Impact}
The \textsf{FeatLLM} framework is crafted to harness the extensive prior knowledge embedded in LLMs, thereby surpassing the limits encountered in conventional supervised tabular learning. By targeting superior data efficiency and leveraging prior knowledge, our approach addresses the prevalent issue of overfitting in contexts with scarce training examples. \textsf{FeatLLM} holds considerable promise for practical domains such as finance and healthcare, where data annotation is both expensive and labor-intensive, and pretrained LLMs may provide vital domain-specific insights. Looking forward, it will be critical to systematically evaluate the effects of societal biases and potential misinformation inherently present in LLMs' prior knowledge, as these may substantially influence model predictions—even overriding limited annotated data. This will be a crucial step for enabling responsible and widespread adoption of the framework.

\end{document}